%-------------------------
% Resume in LaTeX — ATS-optimised
% Based on a template by Aras Gungore (MIT). Modified for: Agam Singh Saini
%
% HOW THIS FILE IS USED
% Edit this file and push to master. GitHub Actions (.github/workflows/deploy-site.yml)
% compiles it with pdflatex and publishes the PDF on the portfolio as
% https://doombuoy.github.io/Portfolio/Agam_Singh_Saini_Resume.pdf
%
% ATS RULES THIS FILE FOLLOWS (keep them when editing)
% - One column, standard section names, reverse-chronological order.
% - No icons or images; links are written out as plain text so parsers can read them.
% - \pdfgentounicode makes every glyph (dashes, ligatures) extract as real text.
% - Dates always "Mon YYYY -- Mon YYYY"; skills written exactly as job ads write them.
% - pdflatex only (no fontspec); write symbols as LaTeX (R$^2$, $\sim$), not Unicode.
%
% HOW THE BULLETS ARE WRITTEN (keep this voice)
% - Lead with what the work made possible, or for whom; the number comes after as proof.
% - Plain verbs (built, cut, found, shipped). No "leveraged", "spearheaded", "robust",
%   "cutting-edge", "seamless", "passionate", "results-driven".
% - Bold only the one result a recruiter should see in each project, not every keyword.
% - Say where a model is weak when it matters. It reads as judgement, not as a flaw.
%
% TAILORING FOR A SPECIFIC JOB (2 minutes)
% 1. Change \targetrole below to the job title in the ad.
% 2. Reorder the Skills lines so the ad's must-haves come first (only list skills you have).
% 3. Use the ad's exact wording for a skill you have (e.g. "Machine Learning" vs "ML").
%------------------------

\documentclass[letterpaper,11pt]{article}

\usepackage[letterpaper, top=0.45in, bottom=0.45in, left=0.5in, right=0.5in]{geometry}
\usepackage[T1]{fontenc}
\usepackage[utf8]{inputenc}
\usepackage[english]{babel}
\usepackage{lmodern}
\usepackage{enumitem}
\usepackage{titlesec}
\usepackage[usenames,dvipsnames]{color}
\usepackage[hidelinks]{hyperref}
\usepackage{fancyhdr}
\usepackage{tabularx}

% Make the PDF text machine-readable (ATS): map every glyph to Unicode
\input{glyphtounicode}
\pdfgentounicode=1

\hypersetup{
  pdftitle={Agam Singh Saini - Data Scientist Resume},
  pdfauthor={Agam Singh Saini},
  pdfsubject={Resume},
  pdfkeywords={Data Scientist, Machine Learning, Python, SQL, PySpark, Databricks, Deep Learning, Graph Neural Networks, LLM, RAG, NLP, Time Series Forecasting, FastAPI, Docker}
}

\newcommand{\targetrole}{Data Scientist \textbar{} Machine Learning Engineer}

\pagestyle{fancy}
\fancyhf{}
\renewcommand{\headrulewidth}{0pt}
\renewcommand{\footrulewidth}{0pt}

\urlstyle{same}
\raggedbottom
\raggedright
\setlength{\tabcolsep}{0in}
\setlist[itemize]{label=-, leftmargin=0.18in, itemsep=0pt, topsep=2pt}

\titleformat{\section}{\vspace{-5pt}\scshape\raggedright\large\bfseries}{}{0em}{}[\titlerule\vspace{-5pt}]

%-------------------------
% Commands
\newcommand{\resumeItem}[1]{\item\small{#1}}

% Role/degree heading: left title, right date; second line organisation, location
\newcommand{\resumeSubheading}[4]{
  \vspace{-1pt}\item
    \begin{tabular*}{0.98\textwidth}[t]{l@{\extracolsep{\fill}}r}
      \textbf{#1} & #2 \\
      \textit{\small #3} & \textit{\small #4} \\
    \end{tabular*}\vspace{-6pt}
}

% Project heading: title + tech on the left, date on the right; links on their own line
\newcommand{\resumeProjectHeading}[3]{
  \vspace{-1pt}\item
    \begin{tabular*}{0.98\textwidth}[t]{l@{\extracolsep{\fill}}r}
      \small #1 & \small #2 \\
    \end{tabular*}\\[-1pt]
    {\footnotesize #3}\vspace{-6pt}
}

\newcommand{\resumeSubHeadingListStart}{\begin{itemize}[leftmargin=0.12in, label={}]}
\newcommand{\resumeSubHeadingListEnd}{\end{itemize}}
\newcommand{\resumeItemListStart}{\begin{itemize}}
\newcommand{\resumeItemListEnd}{\end{itemize}\vspace{-4pt}}

\newcommand{\plainlink}[1]{\href{https://#1}{#1}}

%-------------------------------------------
\begin{document}

%---------- HEADING ----------
\begin{center}
  {\Huge \scshape \textbf{Agam Singh Saini}} \\ \vspace{3pt}
  {\large \targetrole} \\ \vspace{3pt}
  \small
  Sydney, NSW, Australia \textbar{}
  \href{tel:+61425249904}{+61 425 249 904} \textbar{}
  \href{mailto:Agamsingh.saini@student.uts.edu.au}{Agamsingh.saini@student.uts.edu.au} \\ \vspace{1pt}
  \plainlink{linkedin.com/in/agam-singh-saini-199185220} \textbar{}
  \plainlink{github.com/DoomBuoy} \textbar{}
  \plainlink{doombuoy.github.io/Portfolio} \\ \vspace{3pt}
  \textbf{Work rights:} Subclass 500 (48 hrs/fortnight); full working rights (Subclass 485) at time of joining
\end{center}

%----------- SUMMARY -----------
\section{Summary}
\small{
Data scientist finishing a Master of Data Science and Innovation at UTS (December 2026), with a background in mechatronics engineering. My projects tend to end in something people can use: a crime-risk forecast to support police shift planning, a fare model built on a billion New York taxi trips, and a transit app live on the App Store and Google Play. I work in Python and SQL, from cleaning messy data at scale in PySpark to building machine learning models and deploying them with FastAPI and Docker, and I am upfront about where a model should not be trusted.
}

%----------- SKILLS -----------
\section{Technical Skills}
\small{
\begin{itemize}[leftmargin=0.12in, label={}]
  \item \textbf{Programming:} Python, SQL, C++, C, Dart (Flutter)
  \item \textbf{Machine Learning:} Scikit-learn, XGBoost, LightGBM, CatBoost, TensorFlow, deep learning, Graph Neural Networks (GNN), ensemble methods, time series forecasting, classification, regression, computer vision
  \item \textbf{NLP \& Generative AI:} Large Language Models (LLMs), LangGraph, multi-agent systems, Retrieval-Augmented Generation (RAG), Natural Language Inference (NLI), Gemini API, Groq API
  \item \textbf{Data Engineering \& Big Data:} PySpark, Spark SQL, Databricks, Pandas, NumPy, ETL pipelines, data cleaning and validation, Protocol Buffers
  \item \textbf{Deployment \& MLOps:} FastAPI, REST APIs, Streamlit, Docker, Git, GitHub Actions (CI/CD), Azure, Render, Weights \& Biases (experiment tracking)
  \item \textbf{Analytics \& Visualisation:} exploratory data analysis (EDA), Power BI (DAX), Plotly, feature engineering, cross-validation, hyperparameter tuning, model evaluation (F1, ROC-AUC, RMSE, R$^2$), explainable AI (LIME)
\end{itemize}
}
\vspace{-6pt}

%----------- EDUCATION -----------
\section{Education}
\resumeSubHeadingListStart
  \resumeSubheading
    {Master of Data Science and Innovation (GPA 3.5/4)}{Feb 2025 -- Dec 2026}
    {University of Technology Sydney}{Sydney, Australia}
  \resumeSubheading
    {Bachelor of Engineering, Mechatronics (CGPA 8.77/10)}{Mar 2020 -- Jun 2024}
    {Thapar Institute of Engineering and Technology}{Patiala, India}
\resumeSubHeadingListEnd

%----------- PROJECTS -----------
\section{Data Science Projects}
\resumeSubHeadingListStart

\resumeProjectHeading
  {\textbf{CrimeSense: Crime Risk Forecasting with Graph Neural Networks} (MDSI Capstone)}{2026}
  {Python, Graph Neural Networks (STGNN), Gradient Boosting, Streamlit \textbar{} \plainlink{crimes-sense.streamlit.app}}
  \resumeItemListStart
    \resumeItem{In a team of six, built a forecast that gives police shift planners a data-driven starting point instead of last week's incident reports: Low, Moderate or High risk for each of 9,965 US agencies over the next 8 hours, from 33M+ NIBRS records.}
    \resumeItem{Since crime spills across boundaries, linked the agencies into one spatial graph and trained a spatio-temporal GNN that learns from neighbouring areas; it \textbf{beat a Gradient Boosting baseline on every risk class} (macro F1 0.472).}
    \resumeItem{Scored each agency against its own baseline so small towns and big cities are judged fairly, and documented where not to trust the model (only 21\% of High-risk alerts are correct; reported crime reflects where police already patrol).}
  \resumeItemListEnd

\resumeProjectHeading
  {\textbf{Sydly: Real-Time Transit App} (iOS and Android)}{Feb 2026}
  {Flutter, Dart, TfNSW GTFS-R, Protocol Buffers, OpenStreetMap \textbar{} \plainlink{doombuoy.github.io/Transport\_app\_website}}
  \resumeItemListStart
    \resumeItem{Built and \textbf{published a Sydney transit app on the App Store and Google Play} that shows commuters where their bus, train or ferry is right now, with car-park space, cheap fuel and loading zones for drivers in the same app.}
    \resumeItem{Decoded Transport for NSW's binary real-time feeds (GTFS-R, Protocol Buffers) and joined them to timetable data to show 10,000+ moving vehicles on one map, refreshed every 5 seconds.}
  \resumeItemListEnd

\resumeProjectHeading
  {\textbf{DebateCheck: Multi-Agent LLM Health-Claim Checker}}{2026}
  {Python, LangGraph, Gemini, Groq, FastAPI, Streamlit, PubMed API \textbar{} \plainlink{github.com/DoomBuoy/Debate\_check}}
  \resumeItemListStart
    \resumeItem{In a team of five, built a health-claim checker that shows its evidence rather than one unexplained answer: PRO and CON LLM agents (LangGraph) argue over retrieved PubMed abstracts and a separate judge model decides.}
    \resumeItem{It answers ``Unverifiable'' instead of guessing when no study applies, and \textbf{all 181 citations} the agents gave were valid.}
    \resumeItem{Beat an NLI baseline on evidence-backed claims (47.1\% vs 35.3\% accuracy) and showed retrieval was the real bottleneck: PubMed had usable evidence for only 17 of 101 claims.}
  \resumeItemListEnd

\resumeProjectHeading
  {\textbf{NYC Taxi Big Data Analytics and Fare Prediction}}{Nov 2024}
  {PySpark, Spark SQL, Databricks, MLlib \textbar{} \plainlink{github.com/DoomBuoy/New\_York\_Taxi\_Databrics}}
  \resumeItemListStart
    \resumeItem{Built a PySpark ETL pipeline on Databricks that merged a decade of taxi files with shifting schemas into one dataset of about 1 billion trips, far too large for a single machine.}
    \resumeItem{Validation caught 70--80 million bad records (negative fares, impossible distances, broken timestamps) that would have skewed every fare and revenue figure.}
    \resumeItem{The fare model \textbf{predicts 91\% of fares within \$8} on later trips than it was trained on (R$^2$ 0.873; RMSE \$7.63 vs a \$23.05 baseline); Spark SQL analysis found the 6--8 PM peak, tipping habits and the top revenue routes.}
  \resumeItemListEnd

\resumeProjectHeading
  {\textbf{Bitcoin Price Forecasting System}}{Nov 2025}
  {XGBoost, LightGBM, CatBoost, FastAPI, Docker, LIME \textbar{} \plainlink{github.com/DoomBuoy/Digital\_Product\_Bitcoin\_Price\_Prediction}}
  \resumeItemListStart
    \resumeItem{Took a time series model past the notebook into a product: a FastAPI service other apps can call and a public Streamlit dashboard, in Docker, predicting Bitcoin's next-day high from 10 years of data (R$^2$ 0.796).}
    \resumeItem{Added a fallback that keeps the API up if the model fails to load, and LIME explanations for every prediction.}
  \resumeItemListEnd

\resumeProjectHeading
  {\textbf{Rain Forecasting API}}{Nov 2025}
  {CatBoost, XGBoost, LightGBM, FastAPI, Docker \textbar{} \plainlink{github.com/DoomBuoy/Digital\_Product\_Weather\_Prediction}}
  \resumeItemListStart
    \resumeItem{Built week-ahead rain forecasts for people who plan days ahead, like farmers and event organisers: an ensemble classifier for rain in 7 days and a regressor for 3-day rainfall, from 50+ weather features.}
    \resumeItem{Tuned for recall because a missed storm costs more than a false alarm, used time-based cross-validation to avoid leakage, and served both models through a public FastAPI API.}
  \resumeItemListEnd

\resumeProjectHeading
  {\textbf{NBA Draft Prediction}}{2025}
  {Python, Scikit-learn, XGBoost, LDA}
  \resumeItemListStart
    \resumeItem{Predicted each college player's chance of being drafted to show scouts where to look first; a soft-voting ensemble reached \textbf{0.98+ ROC-AUC}, tuned for recall (0.94+) because missing a future star is the costly mistake.}
  \resumeItemListEnd

\resumeProjectHeading
  {\textbf{doombuoy: Open-Source Python Package}}{Nov 2025}
  {Python, Scikit-learn, Pandas \textbar{} \plainlink{github.com/DoomBuoyz/doombuoy}}
  \resumeItemListStart
    \resumeItem{Packaged the model-comparison boilerplate repeated at the start of every project into a TestPyPI library: one call compares models for classification, regression or time series, roughly halving setup code.}
  \resumeItemListEnd

\resumeSubHeadingListEnd

%----------- EXPERIENCE -----------
\section{Experience}
\resumeSubHeadingListStart
  \resumeSubheading
    {Support Worker (Part-Time)}{Apr 2025 -- Present}
    {Informed Care Services}{Sydney, Australia}
    \resumeItemListStart
      \resumeItem{Maintain accurate NDIS care plans and records for 5+ clients and keep families and healthcare providers informed.}
    \resumeItemListEnd

  \resumeSubheading
    {Mechatronics Engineering Intern}{Jun 2023 -- Aug 2023}
    {Spectrum Technology}{Ambala, India}
    \resumeItemListStart
      \resumeItem{Developed LabVIEW IoT data acquisition systems that give operators sub-second readings from 50+ industrial sensors, with automatic alerts and 99.9\% uptime.}
    \resumeItemListEnd

  \resumeSubheading
    {Robotics Engineering R\&D Intern}{Jan 2023 -- Jun 2023}
    {CVRDE, DRDO (Defence Research and Development Organisation)}{Chennai, India}
    \resumeItemListStart
      \resumeItem{Built real-time computer vision pipelines in ROS2 for person detection and obstacle avoidance, improving navigation safety by 40\%; tested over 50+ hours of Gazebo simulation with senior scientists.}
    \resumeItemListEnd
\resumeSubHeadingListEnd

%----------- CERTIFICATIONS -----------
\section{Certifications}
\small{
\begin{itemize}[leftmargin=0.12in, label={}]
  \item DataCamp Data Scientist Professional (2025) \textbar{} IBM Data Science, Coursera (2023) \textbar{} UTS BUILD Global Leadership Program (2025)
  \item NPTEL: Data Analytics with Python; Introduction to Machine Learning (2023) \textbar{} All certificates: \plainlink{github.com/DoomBuoy/Certificates}
\end{itemize}
}
\vspace{-6pt}

%----------- LEADERSHIP -----------
\section{Leadership and Volunteering}
\resumeSubHeadingListStart
  \resumeSubheading
    {Secretary, Chai Collective UTS}{2026 -- Present}
    {University of Technology Sydney}{Sydney, Australia}
    \resumeItemListStart
      \resumeItem{Handle minutes, records and communications, and help run events such as Chai Pe Charcha (101 tickets sold).}
    \resumeItemListEnd
  \resumeSubheading
    {UTS Peer Networker and SOUL Program Volunteer (Graduated)}{Feb 2025 -- 2026}
    {University of Technology Sydney}{Sydney, Australia}
    \resumeItemListStart
      \resumeItem{Helped 1,000+ new students settle in during orientation and supported equity and wellbeing events (about 100 hours).}
    \resumeItemListEnd
  \resumeSubheading
    {Media Ambassador, UTS BUILD}{Nov 2025}
    {University of Technology Sydney}{Sydney, Australia}
  \resumeSubheading
    {Technical Lead, Mechatronics and Robotics Society}{}
    {Thapar Institute of Engineering and Technology}{Patiala, India}
    \resumeItemListStart
      \resumeItem{Ran AI and computer vision workshops for 200+ students; core member of the Adventure Club (20+ expeditions).}
    \resumeItemListEnd
\resumeSubHeadingListEnd

\end{document}
